% Generated draft: requires mathematical and source review.
\documentclass[11pt]{article}
\usepackage[margin=25mm]{geometry}
\usepackage{fontspec}
\setmainfont{Latin Modern Roman}
\usepackage{amsmath,amssymb,mathtools,unicode-math}
\setmathfont{Latin Modern Math}
\usepackage{enumitem,needspace,xurl,hyperref}
\setlength{\parindent}{0pt}
\setlength{\parskip}{5pt}
\setlength{\emergencystretch}{3em}
\newcommand{\N}{\mathbb{N}}
\newcommand{\Z}{\mathbb{Z}}
\newcommand{\Q}{\mathbb{Q}}
\newcommand{\R}{\mathbb{R}}
\newcommand{\C}{\mathbb{C}}
\newcommand{\F}{\mathbb{F}}
\newcommand{\A}{\mathbb{A}}
\newcommand{\PP}{\mathbb{P}}
\newcommand{\E}{\mathbb{E}}
\newcommand{\eps}{\varepsilon}
\newcommand{\dd}{\,\mathrm{d}}
\begin{document}
\section*{UnsolvedMath ID: 30000887}
\textbf{Source problem code:} catalogue entry
\subsection{Definitions and mathematical statement}
For a group \(G\), use the derived series \(G^{(0)}=G\), \(G^{(n+1)}=[G^{(n)},G^{(n)}]\). The first question asks, for every free group \(G\) and every intended index \(n\), whether
\[
 G\cap\bigl(1+\delta_n(\Z G)\bigr)=G^{(n)}.
\]
The intersection uses the standard embedding \(G\subseteq(\Z G)^\times\). The source does not define \(\delta_n\). Its derived-series context suggests a Lie-derived construction in the group ring, but the additive Lie-derived series, the associative ideals it generates, and augmentation-ideal powers are different constructions and cannot be interchanged. The question also requests a description of the analogous intersection over fields of positive characteristic and for arbitrary groups; no particular general formula or equality is supplied for those extensions. The index origin and the precise \(\delta_n\) convention must be recovered before the displayed equality is a self-contained assertion.

Second, let \(G\) be a finite \(p\)-group and \(k=\F_p\). Every \(\sigma\in\operatorname{Aut}(G)\) extends linearly to a \(k\)-algebra automorphism \(\widehat\sigma\) of \(kG\). The requested factorization is
\[
 \forall\alpha\in\operatorname{Aut}_{k\text{-alg}}(kG),\qquad
 \exists\sigma\in\operatorname{Aut}(G)\ \exists\psi\text{ unipotent}:
 \alpha=\psi\circ\widehat\sigma.
\]
This order expresses an automorphism induced by \(G\) followed by a unipotent automorphism. The record does not define unipotent. In the ordinary linear meaning it would require \((\psi-\operatorname{id})^m=0\) on the finite-dimensional vector space \(kG\) for some \(m\); a convention defined through the augmentation filtration would need to be specified separately.

Third, for a field \(K\) of odd positive characteristic, regard \(A=KG\) as a Lie algebra with bracket \([a,b]=ab-ba\). Put \(L^{(0)}=A\) and \(L^{(m+1)}=\operatorname{span}_K\{[a,b]:a,b\in L^{(m)}\}\). Its Lie-derived length \(\operatorname{dl}^L(A)\) is the least \(m\) with \(L^{(m)}=0\). For a group \(B\), \(\operatorname{dl}(B)\) is the least \(m\) with \(B^{(m)}=1\), or infinity if there is none. The question is whether a single function \(f:\N\to\N\), independent of the field, group, and characteristic, satisfies
\[
 \operatorname{dl}(U(KG))\leq f\bigl(\operatorname{dl}^L(KG)\bigr)
\]
for every Lie-solvable \(KG\) in the stated odd-characteristic class. Here \(U(KG)=(KG)^\times\) is the full unit group. No finite-group hypothesis is supplied in this third part, no numerical form of \(f\) is prescribed, and characteristic two is excluded.
\subsection{Short English statement}
The record asks three separate questions about a derived-series intersection in a group ring, a factorization of automorphisms over the prime field, and a universal bound for the derived length of units in odd characteristic. The first two require missing definitions of the ring filtration and of unipotence.
\subsection{Sources}
\begin{itemize}
\item Local UnsolvedMath registry, permanent ID \(30000887\), exact source code \texttt{catalogue entry}, proposed year 2007; no individual proposer is supplied.
\item \emph{Mini-Workshop: Arithmetik von Gruppenringen (2008)}, Oberwolfach Reports, DOI 10.4171/owr/2007/55, pp. 3209--3240. \url{https://doi.org/10.4171/owr/2007/55}.
\item Supplied reference labelled ``derived length of group-algebra units'': \url{https://doi.org/10.21136/CMJ.2017.0205-16}.
\item Supplied reference labelled ``automorphisms of modular group algebras'': \url{https://doi.org/10.1016/j.jalgebra.2023.08.019}.
\item The supplied \texttt{corrected\_verified} note says all three questions and their notation were retained. The assessment checked 2026-08-21 reports important-family results for dimension-subgroup and modular-automorphism cases and an unresolved general derived-length bound. The partially solved label is dated supplied triage requiring independent verification of exact theorem scope.
\end{itemize}
\end{document}
